\chapter{Introduction: A Differentiable Earth System Model}
\label{chap:ch01_introduction}

\section{What is Chronos-ESM?}

\model{} is a coupled Earth System Model (ESM): a numerical representation of the physical
climate system in which an \emph{atmosphere}, an \emph{ocean}, \emph{sea ice}, and a
\emph{land surface} evolve together, exchanging heat, water, and momentum. What distinguishes
it from the large community models (CESM, MPI-ESM, IPSL-CM, and the rest of the CMIP fleet)
is not the physics it contains but \emph{how it is written} and \emph{how fast it runs}:

\begin{itemize}[leftmargin=1.4em]
  \item \textbf{It is differentiable.} The entire model --- every flux, every time step, every
    spectral transform --- is written in \jax{} \citep{jax2018}, a system for composable
    program transformations. As a consequence one can compute the exact gradient of \emph{any}
    scalar diagnostic (global-mean temperature, the strength of the Atlantic overturning, the
    freshwater flux at which that overturning collapses) with respect to \emph{any} input or
    tunable parameter, by automatic differentiation. This turns calibration, sensitivity
    analysis, and variational data assimilation from heroic, model-specific efforts into
    one-line operations (Part~III).
  \item \textbf{It runs on a laptop or a GPU.} At its default T31 spectral resolution
    ($\approx 3.75^\circ$, roughly $96\times48$ horizontal grid points) a coupled century
    integrates in minutes-to-hours on a single GPU, and runs --- more slowly but identically
    --- on a CPU. There is no MPI and no bespoke build system: \code{pip install} and run.
  \item \textbf{It is small enough to understand.} The whole model is about $8{,}000$ lines of
    Python. A student can read the ocean's tracer-advection routine, change the equation of
    state, and watch the climate respond the same afternoon.
\end{itemize}

These three properties --- differentiable, fast, legible --- are what make \model{} suitable
both as a research instrument for the emerging field of differentiable Earth-system science
and as a teaching model for graduate climate courses. This book serves both audiences.

\section{The model at a glance}

\model{} couples four components on a sphere of radius $\Erad$ rotating at rate $\Om$:

\begin{description}[leftmargin=0pt,style=nextline]
  \item[Atmosphere] A multi-level spectral primitive-equation dynamical core built on Google's
    differentiable \dino{} dycore --- the dynamical engine behind NeuralGCM
    \citep{kochkov2024neuralgcm}. It supports genuine baroclinic instability (growing synoptic
    eddies, a mid-latitude jet) and carries moisture for a tropical rain band. A legacy
    single-level spectral model is also provided for the cheapest experiments
    (Chapters~\ref{chap:ch09_atmosphere_multilevel}--\ref{chap:ch10_atmosphere_singlelevel}).
  \item[Ocean] A $z$-level, hydrostatic, Boussinesq primitive-equation ocean with
    Gent--McWilliams eddy transport \citep{gentmcwilliams1990}, initialized from the World
    Ocean Atlas climatology (Chapters~\ref{chap:ch04_ocean_equations}--\ref{chap:ch08_ocean_overturning}).
  \item[Sea ice] A Semtner-style thermodynamic ice model \citep{semtner1976}
    (Chapter~\ref{chap:ch11_seaice}).
  \item[Land] A slab/bucket land-surface scheme with runoff routing
    (Chapter~\ref{chap:ch12_land}).
\end{description}

The components exchange fluxes through a coupler (Chapter~\ref{chap:ch13_coupler_fluxes}) using
bulk aerodynamic formulae, with an optional flux correction that allows the coupled system to
be run either as a tightly observationally-constrained \emph{control} or as a freely-evolving,
forcing-responsive model.

\section{Philosophy: equations and code, side by side}

A recurring frustration in climate science is the distance between the equations a model is
\emph{said} to solve and the equations its code \emph{actually} integrates --- the
discretizations, the limiters, the clamps, the approximations adopted for stability. This book
takes the unusual stance that those details \emph{are} the model and deserve first-class
documentation. Throughout, each governing equation is presented in its continuous form, then in
the discrete form the code uses, with a pointer to the exact source location:

\codref{chronos\_esm/ocean/veros\_driver.py}{the ocean time step \code{step\_ocean}}

\noindent
Where a scheme is a deliberate simplification --- the single-level atmosphere's lack of vertical
structure, or the interim thermohaline closure that stands in for a prognostic overturning until
the momentum core is wired in --- we say so plainly and quantify the consequence. Honest
documentation of an evolving model is itself a scientific contribution.

\section{Why differentiability matters}

Conventional ESMs are evaluated and tuned by running them many times and comparing outputs ---
an expensive, derivative-free search over a high-dimensional parameter space. Because \model{}
is differentiable, the sensitivity of any diagnostic $J$ to any parameter vector
$\boldsymbol{\theta}$ is available directly as $\partial J/\partial\boldsymbol{\theta}$ at the
cost of roughly one extra model evaluation, by reverse-mode automatic differentiation. This
enables:
\begin{itemize}[leftmargin=1.4em]
  \item \emph{gradient-based calibration} --- fit parameters to observations by descent rather
    than by grid search;
  \item \emph{variational data assimilation} (4D-Var) --- find the initial state that best fits
    a window of observations;
  \item \emph{rigorous sensitivity and tipping analysis} --- e.g.\ locating a bifurcation by
    differentiating through the steady-state condition rather than by brute-force sweeps
    (Chapters~\ref{chap:ch15_differentiable_applications} and~\ref{chap:ch17_amoc_tipping}).
\end{itemize}
A theme of Part~III is that differentiability is powerful but not free: some targets --- notably
bifurcations such as an AMOC collapse --- are non-smooth, and require the right formulation
(fold continuation) rather than naive backpropagation through a long integration.

\section{How to read this book}

Part~I develops the continuous and discrete mathematics common to all components, and the
differentiable-programming paradigm. Part~II documents each component model in
equation-complete detail, derived from the code. Part~III covers forcing, sensitivity, and data
assimilation. Part~IV evaluates the model against the canonical phenomena every climate model
must reproduce. Part~V looks outward: a roadmap toward CMIP7 participation, and a guide to using
\model{} in the classroom. Appendices collect the notation, the parameter values, a map from
code modules to chapters, and installation instructions.

Readers wanting a fast orientation may read this chapter, then jump to the test cases of
Part~IV to see what the model produces, returning to Part~II for the mechanisms. Instructors
will find a suggested course outline in Chapter~\ref{chap:ch21_teaching}.

\begin{exercise}
Install \model{} (Appendix~\ref{chap:appD_running}) and run the shortest coupled integration in
the repository. Plot the global-mean surface temperature over the run. How long did it take, and
on what hardware?
\end{exercise}

\begin{exercise}
List three scientific questions you could ask with a differentiable climate model that would be
impractical with a conventional one. For each, identify the scalar diagnostic $J$ and the
parameter $\boldsymbol{\theta}$ whose gradient $\partial J/\partial\boldsymbol{\theta}$ you would
compute.
\end{exercise}
